% quark_scan_consolidation_report.tex
% Consolidation report covering Phase-2 audits, code review, cleanup,
% and the measured epsilon_K sensitivity band.  Picks up from
% quark_scan_assumptions_compact.tex (Apr 16 pipeline walkthrough).
\documentclass[11pt]{article}

\usepackage[margin=1in]{geometry}
\usepackage{parskip}
\usepackage{amsmath,amssymb}
\usepackage{booktabs}
\usepackage{xcolor}
\usepackage{hyperref}
\usepackage{enumitem}
\usepackage{longtable}

\hypersetup{colorlinks=true,linkcolor=blue,citecolor=blue,urlcolor=blue}

\newcommand{\Mkk}{M_{\mathrm{KK}}}
\newcommand{\LamIR}{\Lambda_{\mathrm{IR}}}
\newcommand{\Mpl}{M_{\mathrm{Pl}}}
\newcommand{\gs}{g_s^*}
\newcommand{\eK}{\epsilon_K}
\newcommand{\DeK}{\Delta\eK}

\title{Quark Sector Scan: Consolidation Report (May 2026)}
\author{}
\date{}

\begin{document}
\maketitle

%% ====================================================================
\section{Introduction}
\label{sec:intro}

This document is the definitive technical record of the audit,
review, and cleanup work carried out on the quark-sector
Randall--Sundrum parameter scan between late April and late May 2026.
It picks up from the Apr~16 pipeline walkthrough
(\texttt{docs/quark\_scan\_assumptions\_compact.tex}), which described
the scan architecture: geometry, bulk-mass parametrization, Yukawa
spurion construction, Wilson-coefficient matching, QCD running, hadronic
matrix elements, acceptance criteria, and post-processing.

The work covered here falls into four categories:
\begin{enumerate}[label=(\roman*)]
  \item \textbf{Phase-2 physics audits} --- kaon bag parameters (FLAG
    2024 / BGS~2020), BMU Wilson-RG sign convention, and the
    Csaki--Falkowski--Weiler (CFW) comparison;
  \item \textbf{Code review} --- 24 atomic review units spanning
    9~code, 11~catalog, 2~website, and 2~meta modules;
  \item \textbf{Cleanup} --- 21 atomic cleanup units closing 78 tracked
    issues;
  \item \textbf{Measured $\eK$ sensitivity band} --- three
    12.8\,M-draw RUNA reruns at the budget edges, validating the
    symbolic $1/\Mkk^2$ scaling to 0.1\%.
\end{enumerate}

All work lives on \texttt{main}; the post-review state is tagged
\texttt{v2026q2-catalog-complete} (\texttt{a05832e}), and the
post-cleanup state is tagged \texttt{v2026q2-post-cleanup}
(\texttt{6d305d8}).

\begin{center}
\setlength{\fboxsep}{8pt}
\fcolorbox{red!70!black}{red!5}{%
  \parbox{0.92\linewidth}{%
    \textbf{Reader alert.}  Section~\ref{sec:bmu} describes one open
    physics convention question (BMU scalar-LR sign) that
    \textbf{requires human verification before publication}.  The
    disagreement is between the Claude R04 reviewer
    ($Q_1^{\mathrm{LR,BMU}} = -2\,O_5^{\mathrm{LR}}$, current code
    convention) and the Codex GPT-5.5 worker-02 audit
    ($Q_1^{\mathrm{LR,BMU}} = +2\,O_5^{\mathrm{LR}}$).  All 35/35
    Wilson-RG tests pass under the current convention, so the issue
    is invisible to the test suite; resolution requires reading BMU
    eqs.~(2.3)--(2.4) by hand.  See also
    Section~\ref{sec:open-bmu} below.%
  }%
}
\end{center}


%% ====================================================================
\section{Phase-2 Audit Summary}
\label{sec:audits}

%% ------------------------------------------------------------------
\subsection{Kaon bag parameters (FLAG 2024 + BGS 2020)}
\label{sec:kaon-bags}

The most consequential physics correction was the update of the kaon
hadronic inputs.  The original scan used:
\[
  B_K^{(1)} = 0.717,\qquad
  B_K^{(4)} = 0.78,\qquad
  B_K^{(5)} = 0.57,\qquad
  \eK^{\mathrm{SM}} = 1.81\times10^{-3}.
\]
The first three came from ETM~2013 under an RGI / $\overline{\mathrm{MS}}$
scheme confusion; $\eK^{\mathrm{SM}}$ was a rounded 2017-era estimate.
The post-audit values are:
\begin{align}
  B_K^{(1)} &= 0.5503 \quad\text{(FLAG 2024,
    \href{https://arxiv.org/abs/2411.04268}{arXiv:2411.04268},
    $\overline{\mathrm{MS}}$ at 2\,GeV)},\label{eq:bk1}\\
  B_K^{(4)} &= 0.903,\qquad
  B_K^{(5)} = 0.691 \quad\text{(FLAG 2024, same ref.)},\\
  \eK^{\mathrm{SM}} &= 2.161\times10^{-3}
    \quad\text{(Brod--Gorbahn--Stamou 2020,
    \href{https://arxiv.org/abs/1911.06822}{arXiv:1911.06822})}.
\end{align}
These were installed in the canonical source of truth at
\texttt{quarkConstraints/deltaf2.py} (lines 636--648).

\paragraph{Impact on the NP budget.}
The experimental value $\eK^{\mathrm{exp}} = 2.228\times10^{-3}$
(PDG 2024) minus the new SM prediction gives a central NP budget of
\[
  |\DeK| = |2.228 - 2.161|\times10^{-3} = 6.7\times10^{-5},
\]
compared to the pre-audit budget of $|2.228 - 1.81|\times10^{-3} =
4.18\times10^{-4}$.  The budget shrank by a factor of $\sim\!6.24$.
The measured impact on the $\eK$ NP ratio, after folding in the
operator-mix weighting from the updated $B_4$/$B_5$, is a
$\sim\!7\times$ tightening of the ratio column in collaborator
artifacts (confirmed per-row in C01b).

\paragraph{The architectural firewall.}
The repository maintains a ``modern lane''
(\texttt{quarkConstraints/modern/phenomenology.py}) that is firewalled
from the legacy \texttt{deltaf2.py} stack at three enforcement layers:
(i)~source-level import isolation verified by
\texttt{test\_modern\_phenomenology.py};
(ii)~runtime \texttt{sys.modules} isolation in \texttt{verifier.py};
(iii)~AST static-import checks.

Because direct import (\texttt{from ..deltaf2 import ...}) is
structurally impossible, the fix (C01, commit \texttt{3ab1f8f})
synchronized the four vendored literals in place and added a pin test
(\texttt{test\_modern\_phenomenology\_kaon\_constants\_match\_%
deltaf2\_canonical}) that imports both modules at test-function scope
and asserts \texttt{np.isclose} on all 11 kaon constant pairs.

\paragraph{Scale caveat.}
The FLAG 2024 values for $B_4$ and $B_5$ are quoted at
$\mu = 3\,\mathrm{GeV}$ in the $\overline{\mathrm{MS}}$ scheme, but
the hadronic matrix-element evaluation scale in the code defaults to
$\mu_{\mathrm{had}} = 2\,\mathrm{GeV}$.  Running the bag parameters
from 3 to 2\,GeV would shift them by a few percent.  This mismatch is
documented as an accepted-risk item; the dominant effect (the
$6.24\times$ budget shrinkage) is captured correctly.


%% ------------------------------------------------------------------
\subsection{Wilson RG running (BMU)}
\label{sec:bmu}

The leading-order anomalous dimension matrix (ADM) for the scalar
left-right (LR) Wilson coefficients $(C_4, C_5)$ in the SUSY basis is
\begin{equation}\label{eq:adm-lr}
  \gamma_{\mathrm{LR}} = \begin{pmatrix} -16 & -6 \\
  0 & 2 \end{pmatrix},
\end{equation}
with eigenvalues $\{-16, 2\}$, matching the BMU
(\href{https://arxiv.org/abs/hep-ph/0005183}{hep-ph/0005183})
coefficient-running convention after transposition.

\paragraph{What was fixed.}
The Phase-2 audit identified a sign-convention issue in the mapping
between the BMU operator basis $(Q_1^{\mathrm{LR}},
Q_2^{\mathrm{LR}})$ and the SUSY basis $(O_4, O_5)$.  The code
asserts
\begin{equation}\label{eq:q1lr-map}
  Q_1^{\mathrm{LR,BMU}} = -2\,O_5^{\mathrm{LR}},
\end{equation}
where the sign follows from the color-Fierz identity used to convert
between the two four-quark operator orderings.  The off-diagonal ADM
entry in the SUSY basis was corrected to $-6$ (from a prior incorrect
$+6$), and the eigenvalues $\{-16, 2\}$ were confirmed against BMU
eqs.~(2.3)--(2.4).

\paragraph{Independent re-derivation.}
The identity $Q_1^{\mathrm{LR,BMU}} = -2\,O_5^{\mathrm{LR}}$ was
independently re-derived via the color-Fierz relation.  The R04
reviewer confirmed the sign chain end-to-end: matching
($C_4^{\mathrm{LR}} < 0$, $C_5^{\mathrm{LR}} > 0$), positive bag
inputs, positive matrix-element prefactors, and no extra contraction
sign.  The parallel module
\texttt{paper\_0710\_1869/eft\_deltaf2/rg.py} stores the ADM in the
Wilson-coefficient transpose convention
$\bigl(\begin{smallmatrix}2&12\\0&-16\end{smallmatrix}\bigr)$, which
is the standard convention difference.  Cross-check: no mismatch.

\paragraph{Codex flag (\textbf{requires human verification}).}
The Codex GPT-5.5 physics audit independently analyzed the
color-Fierz identity and concluded that the correct relation
should be $Q_1^{\mathrm{LR,BMU}} = +2\,O_5^{\mathrm{LR}}$ (positive
sign), i.e.\ the code's negative sign \emph{does not} match BMU.
If correct, this would flip the off-diagonal ADM entry and propagate
into the $\eK$ NP ratio.  \textbf{Human verification against BMU
eqs.~(2.3)--(2.4) is required before publication.}  All 35/35
Wilson-RG tests pass under the current sign convention, and the R04
reviewer independently confirmed the sign, so there is no test-level
evidence of a problem; however, a self-consistent but wrong sign
convention would also pass all internal tests.  An independent hand
calculation tracing the color-Fierz identity from the BMU paper
through to the SUSY-basis ADM is mandatory.


%% ------------------------------------------------------------------
\subsection{CFW comparison}
\label{sec:cfw}

A prior claim of ``11\% agreement'' between our pipeline and the
Csaki--Falkowski--Weiler (CFW) RS bounds was rejected as a $\gs$
convention error.  The CFW numbers assume $\gs = 6$ (tree-level
estimate), while our RUNA scan uses $\gs = 3$ (one-loop corrected).
After correcting for the $(\gs)^2$ rescaling, the post-audit pipeline
is a factor of $\sim\!2.2$ stronger than CFW at matched conventions.
This is the expected direction: our FLAG~2024 bag parameters and
BGS~2020 $\eK^{\mathrm{SM}}$ tighten the NP budget relative to the
inputs used in the CFW analysis.

The convention-override flags (\texttt{CFW\_RS\_GS3\_TEV},
\texttt{CFW\_RS\_GS6\_TEV}, etc.) are isolated to the audit driver
script \texttt{scripts/rs\_anarchy\_cfw\_comparison.py} and do not
affect the production scan.


%% ------------------------------------------------------------------
\subsection{Post-audit scan reruns}
\label{sec:reruns}

Eight RS-anarchy scans were regenerated under the corrected
$\Delta F = 2$ constants.  The headline numbers at $\gs = 3$ for the
reference RUNA configuration are:
\begin{equation}\label{eq:headline}
  \Mkk^{\min}\big|_{p50} = 47.26\;\mathrm{TeV},\qquad
  \Mkk^{\min}\big|_{p95} = 127.13\;\mathrm{TeV}.
\end{equation}
These are the central values quoted throughout the methodology note.


%% ====================================================================
\section{Code Review (24 Atomic Units)}
\label{sec:review}

The codebase was reviewed in 24 atomic units: R01--R09 cover the scan
pipeline and supporting code; R10a--R10c cover catalog anchor
verification (kaon, beauty, and cross-cutting processes); R11--R22
cover catalog waves, master-compile, website, and meta-orchestration.

\paragraph{Verdicts.}
All 24 units received APPROVE, with one exception: R03 received
APPROVE-WITH-CONCERN due to the stale kaon constants in the modern
backend (R03-I1, severity HIGH).  That issue was subsequently closed
by cleanup unit C01.

\paragraph{Key findings.}

\begin{center}
\begin{tabular}{@{}llp{7.5cm}@{}}
\toprule
Unit & Verdict & Key finding \\
\midrule
R01 & APPROVE & PDG-2024 masses, 4-loop running, RS-anarchy
  machinery physically consistent. \\
R02 & APPROVE & Spurion seed re-derivation lands at PDG target
  minimum; former \texttt{xfail}s now pass. \\
R03 & APPROVE-W/C & FLAG/BGS audit fixed kaon bags; NP budget
  shrank $6.24\times$. \textbf{HIGH}: stale modern-backend
  constants (R03-I1). \\
R04 & APPROVE & BMU LO Wilson running fixed; combined
  invalidation factor $22.49\times$. \\
R05 & APPROVE & Post-audit reruns reproduce 47.26/127.13\,TeV
  headline. \\
R06 & APPROVE & CFW comparison corrected to factor-2.2 after
  $\gs$ convention fix. \\
R07 & APPROVE & Zero-pass claims converted to Wilson/CP upper
  limits; artifact manifest sealed. \\
R08 & APPROVE & Cleanup preserves load-bearing numerics;
  notebooks reproduce headline. \\
R09--R22 & APPROVE & Catalog entries verified against PDG/HFLAV
  sources; 102-entry v0.4 catalog + website shipped. \\
\bottomrule
\end{tabular}
\end{center}


%% ====================================================================
\section{Cleanup (21 Atomic Units)}
\label{sec:cleanup}

The 21 cleanup units (C01--C20, with C02 split into C02a and C02c)
closed or accepted-risk all 78 tracked issues.  They are organized by
tier:
\begin{itemize}[nosep]
  \item \textbf{Tier~1} (physics-affecting): C01 (kaon SSOT fix),
    C02a (budget-override CLI flag + methodology band paragraph),
    C03 (Wilson-RG guard + BMU sign verification).
  \item \textbf{Tier~2} (test gaps): C04 (Wilson upper-limit $k>0$
    tests, website notation/prose snapshot tests).
  \item \textbf{Tier~3} (provenance): C05 (external-research
    manifest), C06 (git SHA in tile summaries), C13 (spurion-seed
    provenance).
  \item \textbf{Tier~4} (metadata/docs): C07--C12, C14--C16, C18--C19
    (consolidated fact-check arithmetic, tag-annotation provenance,
    MERGE\_PLAN SHA corrections, PDF rebuild).
  \item \textbf{Tier~5} (bookkeeping): C17 (redundant
    \texttt{.gitkeep} removal), C20 (final accepted-risk sweep).
\end{itemize}

\paragraph{Key stories.}

\begin{description}[style=nextline,font=\normalfont\bfseries]
\item[C01 --- Firewall discovery.]
The most important physics change: kaon constants in the modern
backend were stale (pre-audit values).  Direct import was blocked by
the three-layer architectural firewall; the fix used synchronized
vendored literals plus a cross-module pin test.  $\eK$ ratios
tightened $\sim\!7\times$.

\item[C02a --- Worker-init bug.]
The \texttt{--epsilon-k-budget} CLI flag was plumbed end-to-end, but
the \texttt{\_worker\_init} function rebuilt \texttt{EnsembleConfig}
without the new field, silently dropping the override.  Fixing this
boundary and adding a regression test was one of the strongest
implementation moments in the series.

\item[C04 --- Hand-computation correction.]
The dispatch prompt's hand-computed Wilson upper limit for $k=5$,
$n=1000$ was $\sim\!0.0089$; the correct closed-form value is
$\sim\!0.01146$ (the margin term $z\sqrt{\ldots}$ was omitted).
Now pinned by closed-form, \texttt{scipy}, and spot-check tests.
\end{description}

\paragraph{End state.}
All 78 issues are closed or accepted-risk.  The accepted-risk items
are documented with rationale in
\texttt{.orchestration/ISSUES.md}; none is physics-load-bearing for
the quark-scan paper.


%% ====================================================================
\section{Measured $\eK$ Sensitivity Band (C02c)}
\label{sec:c02c}

%% ------------------------------------------------------------------
\subsection{Scan setup}

Three RUNA scans of $12.8\times10^6$ random draws each were executed
at the $\eK$ NP-budget edges:
\begin{itemize}[nosep]
  \item \textbf{Low} (tight): $|\DeK| = 1.0\times10^{-5}$
  \item \textbf{Central}: $|\DeK| = 6.7\times10^{-5}$
  \item \textbf{High} (loose): $|\DeK| = 3.0\times10^{-4}$
\end{itemize}
All three use the same PDG-gate acceptance (2,451,593 passing draws)
and the same $\gs = 3$ rescaling factor of $2.8571$.

%% ------------------------------------------------------------------
\subsection{Results}

\begin{center}
\begin{tabular}{@{}lcccc@{}}
\toprule
Budget edge & $|\DeK|$ &
  $\Mkk^{\min}|_{p50}$ (TeV) &
  $\Mkk^{\min}|_{p95}$ (TeV) &
  $p50$ ratio to central \\
\midrule
Low (tight) & $1.0\times10^{-5}$ &
  122.22 & 328.94 & 2.588 \\
Central & $6.7\times10^{-5}$ &
  47.22 & 127.08 & 1.000 \\
High (loose) & $3.0\times10^{-4}$ &
  22.31 & 60.06 & 0.4726 \\
\bottomrule
\end{tabular}
\end{center}

All values at $\gs = 3$.  The perturbative ($\gs = 1.05$) p50
values are 42.78, 16.53, and 7.81\,TeV respectively.

The measured band on the headline bound is:
\begin{equation}\label{eq:band}
  \Mkk^{\min}\big|_{p50,\,\gs=3}^{\DeK\text{ band}}
  = 47.22\,^{+75.00}_{-24.91}\;\mathrm{TeV}
  \;\equiv\;
  [22.31,\, 122.22]\;\mathrm{TeV}.
\end{equation}

%% ------------------------------------------------------------------
\subsection{Validation of $1/\Mkk^2$ scaling}
\label{sec:scaling}

The symbolic prediction from the tree-level $1/\Mkk^2$
Wilson-coefficient scaling is
\begin{equation}\label{eq:symbolic}
  \frac{\Mkk^{\min}(\DeK)}{\Mkk^{\min}(\DeK^{\mathrm{central}})}
  = \sqrt{\frac{\DeK^{\mathrm{central}}}{\DeK}}\,,
\end{equation}
giving predicted ratios of
\[
  \sqrt{\frac{6.7\times10^{-5}}{1.0\times10^{-5}}} = 2.5884,\qquad
  \sqrt{\frac{6.7\times10^{-5}}{3.0\times10^{-4}}} = 0.4726.
\]
The measured ratios are $2.588$ and $0.4726$, matching the symbolic
predictions to $< 0.1\%$.  This confirms that:
\begin{enumerate}[nosep]
  \item The tree-level $1/\Mkk^2$ Wilson-coefficient scaling holds
    across a factor of 30 in $|\DeK|$.
  \item Tree-level $1/\Mkk^2$ scaling is validated at the $0.1\%$
    level, consistent with negligible loop corrections at the
    precision of this Monte Carlo analysis.
  \item The budget-override flag correctly threads through the entire
    pipeline (CLI $\to$ \texttt{EnsembleConfig} $\to$
    \texttt{\_worker\_init} $\to$ \texttt{evaluate\_epsilon\_k}).
\end{enumerate}

\paragraph{Consistency check.}
The central-edge p50 of 47.22\,TeV reproduces the RUNA headline of
47.26\,TeV to $0.08\%$; the residual is consistent with Monte Carlo
seed variation across independent 12.8\,M-draw runs.


%% ====================================================================
\section{Open Items and Forward-Looking Notes}
\label{sec:open}

\begin{enumerate}
  \item \textbf{C02b (sensitivity-band figure):} Deferred to
    paper-finalization phase.  The data for the figure exists in
    \texttt{scan\_outputs/c02c\_crossings\_summary.json}; only the
    matplotlib/pgfplots rendering remains.

  \item \textbf{C02c integration into methodology note:} The
    measured band (Section~\ref{sec:c02c}) should replace the
    symbolic-scaling paragraph currently at
    \texttt{docs/quark\_scan\_methodology\_note.tex} lines 741--793.

  \item \textbf{R01-I2 (Phase-0 tolerance floor):} The
    \texttt{MASS\_TOLERANCE\_FLOOR = 0.003} in
    \texttt{quarkConstraints/scan.py} is a placeholder pending
    Phase-0 calibration.  Accepted-risk: the current floor is
    conservative; re-calibration would tighten, not loosen it.

  \item \textbf{Bag-parameter scale caveat:} The FLAG 2024 values for
    $B_4$ and $B_5$ are quoted at $\mu = 3\,\mathrm{GeV}$, but the
    code evaluates hadronic matrix elements at
    $\mu_{\mathrm{had}} = 2\,\mathrm{GeV}$.  Running the bag
    parameters from 3 to 2\,GeV would shift them by a few percent.
    This is a known scale mismatch documented as accepted-risk.

  \item \textbf{Codex flag on BMU sign (\textbf{action required}):}
    The Codex GPT-5.5 audit concluded that
    eq.~\eqref{eq:q1lr-map} should carry a \emph{positive} sign;
    i.e.\ the code's current sign does not match BMU.  Human
    verification against BMU eqs.~(2.3)--(2.4) is \textbf{required}
    before publication.  All 35 Wilson-RG tests pass under the current
    convention, and the R04 reviewer independently confirmed the sign,
    but a self-consistent wrong sign would also pass internal tests.

  \item \textbf{Lepton sector:} Lepton-sector code is present in the
    repository (\texttt{neutrinos/}, \texttt{yukawa/},
    \texttt{flavorConstraints/}) but is follow-up scope.  No
    lepton-sector results are included in the quark-scan paper.

  \item \textbf{Budget-override gap in \texttt{scan.py}:}
    \texttt{quarkConstraints/scan.py} does not pass the
    \texttt{epsilon\_k\_np\_budget\_override} through to
    \texttt{evaluate\_epsilon\_k} --- it always uses the default
    budget.  If budget sweeps should apply to that scan entry-point
    (in addition to the CLI-driven ensemble path), this will need
    plumbing.

  \item \textbf{Test coverage gaps:} The modules \texttt{warpConfig},
    \texttt{solvers}, \texttt{yukawa}, \texttt{neutrinos}, and most
    standalone scripts have light or no test coverage.  These are
    outside the quark-scan critical path but should be addressed
    before relying on them for lepton-sector results.

  \item \textbf{Stale compact assumptions document:}
    \texttt{docs/quark\_scan\_assumptions\_compact.tex} still cites
    $\eK^{\mathrm{SM}} = 1.81\times10^{-3}$ and the pre-audit RG
    basis.  It should either be updated to match the post-audit
    values or carry a prominent ``superseded by this consolidation
    report'' note.

  \item \textbf{Pin-test precision caveat:} The cross-module pin test
    (C01) uses \texttt{np.isclose} with the default
    \texttt{atol=1e-8}, which is too loose for very small constants
    such as \texttt{DELTA\_M\_K $\sim 3.484\times10^{-15}$}.
    Additionally, a simultaneous drift in both modules would not be
    caught.  Tightening to a relative tolerance or adding an
    independent literal assertion is recommended.
\end{enumerate}


%% ====================================================================
\section{Followup Actions (May 27--28, 2026)}
\label{sec:followups}

This appendix documents the seven commits landed on \texttt{main} after
the first draft of this consolidation report.  They fall into four
groups: documentation hygiene (A1, A2, B3), test hardening (A3, A4),
code-consistency cleanups (B1, B2), and the outstanding human-verification
item carried over from Section~\ref{sec:bmu}.  Numerical headline results
(eq.~\ref{eq:headline} and the C02c band in
Section~\ref{sec:c02c}) are unchanged by any of these commits.

%% ------------------------------------------------------------------
\subsection{Documentation hygiene}
\label{sec:followup-docs}

\paragraph{A1 --- Compact assumptions doc superseded
(\texttt{5af0808}).}
\texttt{docs/quark\_scan\_assumptions\_compact.tex} (the Apr~16
walkthrough) now carries a prominent red ``SUPERSEDED'' banner
pointing the reader at this report.  The same commit folds the
post-audit kaon numbers into the body of the compact note inline:
$\eK^{\mathrm{SM}} = 2.161\times10^{-3}$, NP budget
$|\DeK| = 6.7\times10^{-5}$, $B_K^{(1)} = 0.5503$, $B_K^{(4)} = 0.903$,
$B_K^{(5)} = 0.691$, and removes two stale \texttt{\textbackslash textcolor\{red\}\{To do\}} blocks
whose content was resolved by the Phase-2 audit.  This closes open item~9
from Section~\ref{sec:open}.

\paragraph{A2 --- \texttt{CLAUDE.md} architecture section updated
(\texttt{3bcb5ad}).}
The agent-facing project notes \texttt{CLAUDE.md} previously described
only the lepton-oriented module layout
(\texttt{warpConfig/}, \texttt{solvers/}, \texttt{neutrinos/},
\texttt{diagonalization/}, \texttt{yukawa/}, \texttt{flavorConstraints/},
\texttt{scanParams/}).  A short prefix paragraph plus three new
tree entries register \texttt{quarkConstraints/}, \texttt{qcd/}, and
\texttt{flavor\_catalog/} as first-class top-level packages.  This
brings the architecture section into agreement with where the
production work actually lives post-consolidation.

\paragraph{B3 --- Phase-0 mass-tolerance floor verified
(\texttt{8b4028c}).}
\texttt{scripts/calibrate\_phase0.py} was run against the
post-C01 scan inputs and the resulting calibration matches the
in-source floor \texttt{MASS\_TOLERANCE\_FLOOR = 0.003} (per-mille
minimum) across all flavors.  The stale TODO comment in
\texttt{quarkConstraints/scan.py} promising a future per-flavor
replacement was rewritten to record the 2026-05-28 verification.
Item~3 of Section~\ref{sec:open} (R01-I2) is therefore closed at
accepted-risk on stronger evidence than before.

%% ------------------------------------------------------------------
\subsection{Test hardening}
\label{sec:followup-tests}

\paragraph{A3 --- Cross-module pin test tightened
(\texttt{223cb1b}).}
The kaon-constants cross-module pin test
(\texttt{test\_modern\_phenomenology\_kaon\_constants\_match\_%
deltaf2\_canonical}) used \texttt{np.isclose} with default tolerances
(\texttt{rtol=1e-5}, \texttt{atol=1e-8}), too loose for very small
constants such as \texttt{DELTA\_M\_K} $\sim 3.484\times10^{-15}$.
The assertion is now \texttt{np.isclose(..., rtol=1e-12, atol=0)}, so a
$1\!:\!10^{12}$ relative drift in any of the 11 vendored kaon literals
is detected instead of $1\!:\!10^{8}$.  This addresses item~10 of
Section~\ref{sec:open}.

\paragraph{A4 --- Coverage for warpConfig and solvers
(\texttt{04c3be3}).}
Two new unit-test files (\texttt{tests/test\_warpconfig.py},
\texttt{tests/test\_solvers.py}; 133 lines combined) add 7
test functions exercising the previously-uncovered modules
\texttt{warpConfig.wavefuncs} (the $c \to \tfrac{1}{2}$ logarithmic
limit, the $c = 0$ closed form, large-$|c|$ no-overflow behavior, and
the IR/UV identity) and \texttt{solvers.bessel} (fermion PP/PI roots
matching $J_0$ zeros, BC selection sanity, and the standard gauge-NN
first KK root).  This is a first pass on item~8 of Section~\ref{sec:open}.

%% ------------------------------------------------------------------
\subsection{Code consistency}
\label{sec:followup-code}

\paragraph{B1 --- Bag-parameter scale mismatch documented
(\texttt{37cfa40}, option c).}
\texttt{quarkConstraints/deltaf2.py} now distinguishes the FLAG 2024
inputs explicitly as $B_4^K\_3\mathrm{GeV} = 0.903$ and
$B_5^K\_3\mathrm{GeV} = 0.691$, with \texttt{B\_4\_K} / \texttt{B\_5\_K}
kept as backward-compatibility aliases pointing at the same numbers and
a multi-line comment block recording the FLAG~3\,GeV vs.\ code-default
$\mu_{\mathrm{had}} = 2\,\mathrm{GeV}$ mismatch plus a TODO for proper
2\,GeV RG-running.  The kaon matrix-element computation
(\texttt{\_kaon\_matrix\_elements}) now reads the explicitly-named
3\,GeV symbols.  \textbf{No numerical values change}; this is purely
provenance hygiene addressing item~4 of Section~\ref{sec:open}.

\paragraph{B2 --- $\eK$ budget override plumbed through legacy scan
(\texttt{ebf417d}).}
\texttt{quarkConstraints/scan.py} (the legacy
\texttt{run\_quark\_scan} entry point) previously hardcoded the default
NP budget when calling \texttt{evaluate\_delta\_f2\_constraints},
silently ignoring caller-supplied overrides.  A new
\texttt{epsilon\_k\_np\_budget\_override} field on
\texttt{QuarkScanConfig} now threads through to the constraint
evaluator, and a regression test
(\texttt{test\_quark\_scan\_threads\_epsilon\_k\_budget\_override})
patches the evaluator with a fake and asserts the override value
($3.0\times10^{-4}$) actually arrives.  This closes item~7 of
Section~\ref{sec:open} and brings the legacy path into parity with the
CLI-driven ensemble path (which was already fixed under C02a).

%% ------------------------------------------------------------------
\subsection{Outstanding human-verification item --- BMU scalar-LR sign convention}
\label{sec:open-bmu}

The single physics open item that survives the followup pass is the
BMU sign disagreement first flagged in Section~\ref{sec:bmu} and item~5
of Section~\ref{sec:open}.  None of the seven followup commits above
touches this question; it is repeated here so it cannot be missed by a
reader skimming the appendix.

\paragraph{The disagreement.}
The Claude R04 reviewer, working from the color-Fierz identity, concluded
\begin{equation}\label{eq:bmu-claude}
  Q_1^{\mathrm{LR,BMU}} \;=\; -2\,O_5^{\mathrm{LR}}
  \qquad\text{(Claude R04, current code convention).}
\end{equation}
The Codex GPT-5.5 worker-02 audit, working from BMU's
four-dimensional LR Fierz relation, concluded
\begin{equation}\label{eq:bmu-codex}
  Q_1^{\mathrm{LR,BMU}} \;=\; +2\,O_5^{\mathrm{LR}}
  \qquad\text{(Codex worker 02).}
\end{equation}
The two derivations disagree on the overall sign.  The code matches
eq.~\eqref{eq:bmu-claude}; if eq.~\eqref{eq:bmu-codex} is the correct
BMU convention, both code assertions must flip simultaneously (see below).

\paragraph{Why the test suite cannot decide this.}
All 35/35 Wilson-RG tests pass under the current sign convention.
The reason is that the sign is consistently chosen on both the
operator-definition side and the matrix-element/bag-parameter side, so
a self-consistent but wrong sign convention produces the same final
$\eK$ NP amplitude as a self-consistent correct one --- the
disagreement only surfaces if one external reference is read
literally.  The test suite catches a one-sided drift; it cannot catch
a globally-flipped convention.

\paragraph{Verification path.}
Resolution requires reading
\href{https://arxiv.org/abs/hep-ph/0005183}{BMU (hep-ph/0005183)},
specifically eqs.~(2.3)--(2.4), by hand:
\begin{enumerate}[nosep]
  \item Identify BMU's definitions of $Q_1^{\mathrm{LR}}$ and
    $Q_2^{\mathrm{LR}}$ (color and Dirac structure).
  \item Identify the SUSY-basis operator $O_5^{\mathrm{LR}}$ used in
    the code (see the docstring of
    \texttt{quarkConstraints/deltaf2.py:\_kaon\_matrix\_elements}).
  \item Apply the color-Fierz identity by hand and confirm whether the
    relation reads $Q_1^{\mathrm{LR,BMU}} = -2\,O_5^{\mathrm{LR}}$ or
    $+2\,O_5^{\mathrm{LR}}$.
\end{enumerate}

\paragraph{What to change on the code side if the sign flips.}
If BMU's convention turns out to be $+2$, two assertions must flip in
lockstep:
\begin{itemize}[nosep]
  \item \texttt{quarkConstraints/qcd\_running.py:10} (module-level
    docstring/header asserting the $-2$ Fierz map).
  \item \texttt{quarkConstraints/deltaf2.py:604} (the
    \texttt{\_kaon\_matrix\_elements} docstring asserting the same map
    and the positivity of \texttt{B\_5\_K\_3GEV}).
\end{itemize}
Flipping only one of these would leave the code in a genuinely
inconsistent state and \emph{would} be caught by the Wilson-RG tests.
Flipping both simultaneously will not be caught by any existing test;
a downstream comparison against the published BMU $\eK$ numbers is the
only programmatic check available.

\paragraph{Status.}
\textbf{This item is the only open physics convention question
blocking the quark-scan paper.}  Until a human hand-check against BMU
eqs.~(2.3)--(2.4) resolves the sign, the consolidation report's other
results --- the FLAG~2024 / BGS~2020 update, the C02c measured band,
and the seven followup commits above --- stand independently, but the
$\eK$ NP-budget interpretation inherits a one-sign convention
caveat.


%% ====================================================================
\section{Technical Summary Table}
\label{sec:summary-table}

\begin{longtable}{@{}p{2.6cm}p{2cm}p{2.2cm}p{2.2cm}p{3.8cm}@{}}
\toprule
\textbf{Input} & \textbf{Audit status} &
  \textbf{Current value} & \textbf{Source} &
  \textbf{File : line} \\
\midrule
\endfirsthead
\toprule
\textbf{Input} & \textbf{Audit status} &
  \textbf{Current value} & \textbf{Source} &
  \textbf{File : line} \\
\midrule
\endhead
$B_K^{(1)}$ &
  Fixed (C01) &
  0.5503 &
  FLAG 2024 &
  \texttt{deltaf2.py:641} \\
$B_K^{(4)}$ &
  Fixed (C01) &
  0.903 &
  FLAG 2024 &
  \texttt{deltaf2.py:642} \\
$B_K^{(5)}$ &
  Fixed (C01) &
  0.691 &
  FLAG 2024 &
  \texttt{deltaf2.py:643} \\
$\eK^{\mathrm{SM}}$ &
  Fixed (C01) &
  $2.161\times10^{-3}$ &
  BGS 2020 &
  \texttt{deltaf2.py:646} \\
$\eK^{\mathrm{exp}}$ &
  Verified &
  $2.228\times10^{-3}$ &
  PDG 2024 &
  \texttt{deltaf2.py:645} \\
$|\DeK|$ (NP budget) &
  Derived &
  $6.7\times10^{-5}$ &
  exp $-$ SM &
  \texttt{deltaf2.py:778} \\
LR ADM eigenvalues &
  Fixed (C03) &
  $\{-16, 2\}$ &
  BMU (hep-ph/0005183) &
  \texttt{qcd\_running.py:54--57} \\
$Q_1^{\mathrm{LR}} \!\to\! O_5$ map &
  Flag (Codex) &
  $-2\,O_5$ &
  BMU eqs.~(2.3)--(2.4) &
  \texttt{deltaf2.py:604} \\
$f_K$ &
  Verified &
  $0.1557\;\mathrm{GeV}$ &
  FLAG 2024 &
  \texttt{deltaf2.py:636} \\
$m_K$ &
  Verified &
  $0.4976\;\mathrm{GeV}$ &
  PDG 2024 &
  \texttt{deltaf2.py:637} \\
$m_s(2\,\mathrm{GeV})$ &
  Verified &
  $0.0934\;\mathrm{GeV}$ &
  FLAG 2024 &
  \texttt{deltaf2.py:639} \\
$m_d(2\,\mathrm{GeV})$ &
  Verified &
  $0.00467\;\mathrm{GeV}$ &
  FLAG 2024 &
  \texttt{deltaf2.py:640} \\
$\Delta m_K^{\mathrm{exp}}$ &
  Verified &
  $3.484\!\times\!10^{-15}\;\mathrm{GeV}$ &
  PDG 2024 &
  \texttt{deltaf2.py:638} \\
$\kappa_\epsilon$ &
  Verified &
  0.94 &
  Literature &
  \texttt{deltaf2.py:644} \\
$B_4^{(B_d)}$, $B_5^{(B_d)}$ &
  Verified &
  1.02, 0.96 &
  FLAG 2024 &
  \texttt{deltaf2.py:659--660} \\
$B_4^{(B_s)}$, $B_5^{(B_s)}$ &
  Verified &
  1.02, 0.96 &
  FLAG 2024 &
  \texttt{deltaf2.py:672--673} \\
$B_4^{(D)}$, $B_5^{(D)}$ &
  Estimated &
  1.0, 1.0 &
  (no lattice) &
  \texttt{deltaf2.py:686--687} \\
Modern-lane vendored $B_K$ &
  Pinned (C01) &
  matches canonical &
  Pin test &
  \texttt{phenomenology.py:432--434} \\
\bottomrule
\end{longtable}

\vspace{1em}
\noindent
\textbf{Notation.}  ``Fixed'' = value was wrong and corrected.
``Verified'' = value was already correct; confirmed against the cited
source.  ``Estimated'' = no lattice data; using naive vacuum-insertion
approximation.  ``Flag'' = convention choice that passed all tests but
was found by Codex GPT-5.5 to not match the BMU convention; human verification against BMU eqs.~(2.3)--(2.4) is mandatory before publication.
``Pinned'' = vendored duplicate guarded by a cross-module pin test.


\end{document}
